
% ===========================================================================
% 2. PROBLEMATIQUE ET SOLUTION
% ===========================================================================
\section{Probl\'ematique et solution}

\subsection{Le constat (probl\`eme)}

Les \'equipes en charge de la s\'ecurit\'e des syst\`emes d'information font
face \`a plusieurs difficult\'es r\'ecurrentes :

\begin{itemize}[leftmargin=1.6em]
  \item \textbf{Scans manuels et ponctuels} : les audits de vuln\'erabilit\'es
        sont r\'ealis\'es de mani\`ere ponctuelle et artisanale, avec des
        outils disparates et des p\'erim\`etres partiels (liste fixe de ports,
        fr\'equence irr\'eguli\`ere).
  \item \textbf{Absence de vision centralis\'ee} : les r\'esultats sont
        \'eparpill\'es dans des fichiers et aucun indicateur global n'est
        disponible pour piloter la s\'ecurit\'e de l'infrastructure.
  \item \textbf{Pas de corr\'elation avec les menaces} : une faille ancienne
        \`a fort impact est trait\'ee de la m\^eme mani\`ere qu'une faille
        r\'ecemment exploit\'ee dans la nature ; aucun contexte de menace
        (EPSS, catalogue KEV, flux d'intelligence MISP) n'est pris en compte.
  \item \textbf{Notation de risque incoh\'erente} : la seule note CVSS ne
        refl\`ete pas l'exposition r\'eelle de l'entreprise (actif critique,
        exploit public, exposition r\'eseau, etc.).
  \item \textbf{Suivi de correction manuel} : l'ouverture et le suivi des
        tickets de rem\'ediation sont effectu\'es \`a la main, ce qui
        g\'en\`ere des oublis et des d\'elais.
  \item \textbf{Observabilit\'e h\'et\'erog\`ene} : les journaux, les
        m\'etriques et la disponibilit\'e des services sont suivis sans outil
        commun ni historique centralis\'e.
\end{itemize}

\subsection{La solution apport\'ee}

La \textbf{VOC Platform} automatise l'ensemble de ce cycle de vie :

\begin{itemize}[leftmargin=1.6em]
  \item \textbf{Scan automatique et complet} : d\'ecouverte du r\'eseau
        (\texttt{nmap -sn}), scan de \textbf{tous les ports TCP} (65535 ports,
        \texttt{-p-}), d\'etection des versions de services (\texttt{-sV}) et
        \textbf{empreinte du syst\`eme d'exploitation} (\texttt{-O}).
  \item \textbf{Analyse diff\'erentielle} : chaque r\'esultat est compar\'e
        \`a l'\'etat pr\'ec\'edent (m\'emoris\'e dans Redis) pour distinguer
        les vuln\'erabilit\'es \emph{nouvelles}, \emph{toujours pr\'esentes}
        ou \emph{r\'esolues}.
  \item \textbf{Enrichissement et corr\'elation} : corr\'elation avec le flux
        CVE de la NVD, recherche d'IOC dans MISP, ajout du score EPSS, de la
        pr\'esence dans le catalogue CISA KEV, d'Exploit-DB et d'OSV.
  \item \textbf{Score de risque m\'etier} : un moteur de risque d\'edi\'e
        (FastAPI) combine CVSS, criticit\'e de l'actif, activit\'e de menace,
        exploitabilit\'e, exposition r\'eseau, EPSS et KEV pour produire une
        note normalis\'ee entre 0 et 10.
  \item \textbf{Visualisation centralis\'ee} : tous les r\'esultats sont
        index\'es dans Elasticsearch et consultables dans Kibana via des
        tableaux de bord.
  \item \textbf{Circuit ITSM automatis\'e} : toute vuln\'erabilit\'e dont le
        risque est sup\'erieur ou \'egal \`a 7 g\'en\`ere automatiquement un
        ticket dans GLPI.
  \item \textbf{Partage de renseignements} : les nouvelles vuln\'erabilit\'es
        sont publi\'ees sous forme d'\'ev\'enements MISP.
  \item \textbf{Observabilit\'e unifi\'ee} : les agents Elastic (Fleet et
        Beats) collectent les journaux, les m\'etriques, les flux r\'eseau,
        les traces d'audit et la disponibilit\'e des services.
\end{itemize}

\subsection{Objectifs du projet}

Les objectifs fix\'es au d\'ebut du stage \'etaient les suivants :

\begin{enumerate}[leftmargin=1.6em]
  \item d\'eployer une plateforme de scan automatique fiable et compl\`ete ;
  \item assurer la corr\'elation des d\'ecouvertes avec les sources de
        menaces ;
  \item produire un score de risque orient\'e m\'etier, reproductible et
        explicable ;
  \item offrir une visibilit\'e centralis\'ee (tableaux de bord) ;
  \item automatiser l'ouverture et le suivi des tickets de rem\'ediation ;
  \item superviser la plateforme et les agents d\'eploy\'es.
\end{enumerate}

% ===========================================================================
% 3. ENVIRONNEMENT DE STAGE
% ===========================================================================
\section{Environnement de stage}

\subsection{Universit\'e TEK-UP University}

TEK-UP University est un \'etablissement d'enseignement sup\'erieur
d\'elivrant des formations d'ing\'enierie dans les domaines des technologies
de l'information et de la communication. Le pr\'esent stage s'inscrit dans le
cadre de l'ann\'ee universitaire 2025--2026 et du programme de formation
d'ing\'enierie.

\subsection{Soci\'et\'e G\'erence Informatique}

G\'erence Informatique est une soci\'et\'e sp\'ecialis\'ee dans les services
li\'es aux syst\`emes d'information : administration de r\'eseaux et de
serveurs, gestion des parcs informatiques et s\'ecurit\'e des syst\`emes.
C'est dans ce contexte que le besoin d'une plateforme centralis\'ee de gestion
des vuln\'erabilit\'es a \'et\'e identifi\'e, afin de professionnaliser la
surveillance s\'ecuritaire de l'infrastructure.

\subsection{Encadrement}

Le stage a \'et\'e encadr\'e par \textbf{Mr.~MDAOUKHI Adel}, qui a assur\'e le
suivi technique du projet, la validation des choix d'architecture et la revue
des d\'eveloppements r\'ealis\'es.

\subsection{M\'ethodologie suivie}

Le travail a \'et\'e conduit selon une d\'emarche it\'erative :

\begin{enumerate}[leftmargin=1.6em]
  \item \textbf{Analyse des besoins} : recensement des difficult\'es, des
        p\'erim\`etres r\'eseau et des outils d\'ej\`a pr\'esents sur
        l'infrastructure ;
  \item \textbf{Conception} : choix de la stack technique, d\'efinition de
        l'architecture et du pipeline de traitement ;
  \item \textbf{D\'eveloppement} : r\'ealisation des modules (scan,
        enrichissement, score, indexation, tickets), des diagrammes et des
        tests unitaires ;
  \item \textbf{D\'eploiement et int\'egration} : mise en place des services
        Docker, configuration des agents Elastic, r\'esolution des
        probl\`emes d'exploitation ;
  \item \textbf{Validation et exploitation} : v\'erification de bout en bout,
        analyse des r\'esultats et suivi op\'erationnel.
\end{enumerate}
